\section{Priority function}
\label{sec:priority}

\begin{lede}
Who plans first each timestep. Robots are ranked by what they are carrying, and
a robot that has been waiting long enough always eventually outranks everyone
--- that bound is the fairness guarantee.
\end{lede}

The priority function must satisfy two competing goals: a loaded robot should
almost always outrank a free one, so deliveries do not stall behind idle
repositioning, and yet \textbf{no robot may starve forever}. LDA-PIBT gets both
from a class term plus an unbounded linear-in-waiting term:
\begin{equation}
  p_i(t) = P_{\mathrm{class}}(i)
    + k_w W_i + k_b B_i + k_u U_i + k_e E_i + \varepsilon_i,
  \label{eq:priority}
\end{equation}

\begin{center}\small
\begin{tabular}{@{}lll@{}}
\toprule
\textbf{Term} & \textbf{Meaning} & \textbf{Parameter} \\
\midrule
$P_{\mathrm{class}}(i)$ & task-class base priority (below) & \code{task\_class\_priority} \\
$k_w W_i$ & waiting weight $\times$ \code{robot.waiting\_time} & \param{waiting\_weight} $= 5$ \\
$k_b B_i$ & blocked weight $\times$ \code{robot.blocked\_time} & \param{blocked\_weight} $= 10$ \\
$k_u U_i$ & urgency weight $\times$ task urgency (0 if no task) & \param{urgency\_weight} $= 1$ \\
$k_e E_i$ & bonus if the robot is currently inside an aisle & \param{priority\_inside\_aisle} $= 50$ \\
$\varepsilon_i$ & tie-breaker, $10^{-4} \cdot \mathrm{id}$ & --- \\
\bottomrule
\end{tabular}
\end{center}

\subsection{Task-class priority}
\label{sec:priority-class}

\begin{equation}
  P_{\mathrm{class}}(i) =
  \begin{cases}
    400 & \text{state} = \textsc{recovery},\\
    300 & \text{task status} = \textsc{to\_delivery},\\
    200 & \text{task status} = \textsc{to\_pickup},\\
    100 & \text{no task, state} = \textsc{parked},\\
    0 & \text{otherwise.}
  \end{cases}
\end{equation}

The ordering $P_{\mathrm{emergency}} > P_{\mathrm{loaded}} >
P_{\mathrm{pickup}} > P_{\mathrm{repositioning}} > P_{\mathrm{free}}$ says: a
robot mid-recovery outranks everything; a robot carrying a delivery outranks
one still travelling to a pickup; an idle robot heading for a parking bay
slightly outranks a plain idle one, so it is not shoved off its chosen bay by
another idle robot.

\subsection{Fairness horizon}
\label{sec:priority-fairness}

Because $k_w W_i$ grows without bound while $P_{\mathrm{class}}$ is fixed, any
robot eventually outranks any other, however low its class. The time this takes
is exact:
\begin{equation}
  \mathrm{fairness\_horizon}
    = \frac{P_{\mathrm{emergency}} - P_{\mathrm{free}}}{k_w}
    = \frac{400 - 0}{5.0} = 80 \ \text{timesteps.}
  \label{eq:fairness-horizon}
\end{equation}

After 80 steps of waiting, a $P_{\mathrm{free}}$ robot's waiting term alone
exceeds the entire class spread, and it outranks even an emergency robot with
$W = 0$. This is the guarantee \code{test\_lifelong.py} checks directly. If
$k_w \leq 0$ the function returns $+\INF$ --- fairness would no longer be
guaranteed, and it surfaces that rather than silently allowing it.

% worked-example: priority
\begin{workedexample}
\textbf{Worked example.} The same scenario run out to timestep 149, so that robots have had time to accumulate waiting. Three of the 24: the highest-priority robot, the one that has waited longest, and the lowest-priority one.

{\fontsize{8.5}{10.2}\selectfont
\begin{verbatim}
robot     state        task         P_class  k_w·W  k_b·B  k_e·E  p_i(t)
--------  -----------  -----------  -------  -----  -----  -----  -------
robot 15  to_delivery  to_delivery  300      0      0      50     350.002
robot 3   to_pickup    to_pickup    200      15     0      50     265
robot 0   to_pickup    to_pickup    200      0      0      50     250
\end{verbatim}
}

\code{P\_class} sets the band; the waiting and blocked terms move a robot within it and, given long enough, out of it. The full class spread is 400 (400 for a robot in recovery down to 0 for an idle one) and \code{k\_w} = 5, so 80 steps of waiting is worth the entire spread: a robot that waits that long outranks \emph{any} robot of any class (\Cref{sec:priority-fairness}).

That bound is the fairness guarantee, and it is the reason the priority function has a waiting term at all. Without it a permanently higher-class neighbour — a loaded robot on a floor that always has loaded robots — could starve an idle one indefinitely, and the throughput number would never show it.
\end{workedexample}
% /worked-example

\subsection{Tie-breaking and ordering}

\code{order\_by\_priority} sorts by $(-p_i, \mathrm{id})$: descending priority,
id as a deterministic secondary key. Combined with the $\varepsilon_i$ term
baked into $p_i$ itself, ties are essentially never reached in practice, but
the explicit id key guarantees determinism regardless.
